\chapter*{Introducción}
\addcontentsline{toc}{chapter}{Introducción}

Diseñar y construir un sistema operativo desde cero es una tarea que exige comprender a fondo las distintas maneras en que este tipo de software puede organizarse y funcionar. Más allá de la teoría, resulta especialmente útil examinar sistemas operativos reales —muchos de ellos abiertos y creados con fines educativos o de investigación— que exploran diferentes arquitecturas, modelos de aislamiento y niveles de complejidad.

En este informe se estudian diversas propuestas existentes, entre ellas MINIX, XV6, Theseus, RedoxOS, FlexOS, LibrettOS, FreeRTOS, HelenOS, Haiku OS y FreeDOS. El propósito es extraer lecciones de diseño, comparar enfoques y evaluar cuáles de estos proyectos pueden servir mejor como referencia para el desarrollo de un sistema operativo propio dentro del ámbito académico.

\section*{Objetivo de la investigación}

El objetivo central del trabajo es ofrecer un análisis técnico, comparativo y crítico de varios sistemas operativos representativos de distintos modelos arquitectónicos y enfoques de diseño, poniendo especial atención en aquellos que son de código abierto y se emplean con fines educativos o de investigación.

\section*{Alcance del análisis}

Este estudio se limita a sistemas operativos que, por su licencia, documentación y propósito, pueden emplearse adecuadamente como referencia en la construcción de un sistema operativo académico. En particular, se revisan los siguientes proyectos: MINIX, XV6, Theseus, RedoxOS, FlexOS, LibrettOS, FreeRTOS, HelenOS, Haiku OS y FreeDOS, cada uno abordado en el Capítulo 1 bajo un esquema de análisis uniforme.

El alcance comprende únicamente:

\begin{enumerate}
\item La revisión de documentación oficial, artículos académicos y repositorios de código.
\item El estudio cualitativo de las decisiones de diseño adoptadas por cada sistema.
\item La elaboración de una comparación estructurada y un análisis crítico orientado a identificar qué propuestas resultan más útiles como inspiración para un nuevo sistema operativo académico.
\end{enumerate}

\section*{Metodología}

Para cumplir con los objetivos planteados, se adopta una metodología dividida en tres fases principales:

\begin{enumerate}
\item \textbf{Revisión bibliográfica y documental:} Se recopila y analiza información obtenida de libros especializados, artículos científicos, documentación oficial de los proyectos y repositorios de código fuente. Esta etapa permite comprender el contexto, los objetivos originales y el estado actual de cada sistema operativo estudiado.

\item \textbf{Evaluación técnica de cada sistema operativo:} En el Capítulo 1 se examina cada propuesta siguiendo una plantilla común que incluye: nombre y propósito del proyecto, arquitectura general, componentes implementados (procesos, memoria, archivos, dispositivos, etc.), herramientas utilizadas (compiladores, emuladores, toolchains) y nivel de complejidad y accesibilidad para estudiantes.

\item \textbf{Comparación estructurada y análisis:} A partir de la información reunida, el Capítulo 2 presenta una comparación global entre los sistemas, considerando criterios como la arquitectura, el tamaño del kernel, el soporte de hardware, la comunidad, la documentación y su utilidad docente. Por último, en el Capítulo 3 se desarrolla un análisis crítico en el que se discuten las ventajas y limitaciones de cada propuesta, reflexionando sobre cuáles son más adecuadas como punto de partida o referencia para la construcción de un sistema operativo desde cero.
\end{enumerate}

Esta metodología permite abordar el estudio de forma ordenada y sistemática, asegurando una comprensión sólida y una evaluación crítica de las distintas propuestas en el contexto académico.